\begin{figure}[htbp]
\centering
\resizebox{0.95\linewidth}{!}{%
\begin{tikzpicture}[
    node distance = 0.85cm and 0cm,
    box/.style = {
        rectangle, 
        rounded corners=4pt, 
        draw=blue!70!black, 
        fill=blue!4, 
        line width=0.8pt, 
        text width=14.2cm, 
        align=center,
        inner sep=6pt,
        font=\small
    },
    arrow/.style = {
        -{Stealth[scale=1.1]}, 
        line width=1.1pt, 
        draw=blue!60!black
    }
]

% Node 1
\node (n1) [box] {
    \textbf{Giai đoạn 1945 -- 1954: Kháng chiến kiến quốc và đặt nền móng}\\
    \textit{Cơ quan chủ quản:} Nha Kỹ thuật vụ (Bộ Quốc gia Giáo dục, theo Sắc lệnh số 119/SL)\\
    \textit{Mô hình:} Đào tạo công nhân kỹ thuật quốc phòng, phục vụ công xưởng chiến khu
};

% Node 2
\node (n2) [box, below=of n1] {
    \textbf{Giai đoạn 1954 -- 1975: Xây dựng chủ nghĩa xã hội ở miền Bắc và chi viện miền Nam}\\
    \textit{Cơ quan chủ quản:} Vụ Quản lý nhân công $\rightarrow$ Tổng cục Đào tạo công nhân kỹ thuật (1969)\\
    \textit{Mô hình:} Đào tạo theo ngành công nghiệp nặng; hợp tác với Liên Xô và Đông Âu
};

% Node 3
\node (n3) [box, below=of n2] {
    \textbf{Giai đoạn 1975 -- 1986: Thống nhất đất nước và kinh tế kế hoạch hóa bao cấp}\\
    \textit{Cơ quan chủ quản:} Tổng cục Dạy nghề trực thuộc Hội đồng Chính phủ (Nghị định số 240/CP)\\
    \textit{Mô hình:} Đào tạo theo chỉ tiêu pháp lệnh; Nhà nước bao cấp và phân bổ việc làm
};

% Node 4
\node (n4) [box, below=of n3] {
    \textbf{Giai đoạn 1986 -- 2006: Bước ngoặt Đổi mới và chuyển giao thể chế}\\
    \textit{Cơ quan chủ quản:} Bộ Giáo dục và Đào tạo $\rightarrow$ Tổng cục Dạy nghề thuộc Bộ Lao động (1998)\\
    \textit{Mô hình:} Đa dạng hóa loại hình cơ sở dạy nghề thích ứng kinh tế thị trường (Quyết định 253)
};

% Node 5
\node (n5) [box, below=of n4] {
    \textbf{Giai đoạn 2006 -- 2014: Thể chế hóa chuyên biệt theo Luật Dạy nghề 2006}\\
    \textit{Cơ quan chủ quản:} Tổng cục Dạy nghề (Bộ Lao động -- Thương binh và Xã hội)\\
    \textit{Mô hình:} Ba cấp trình độ (Sơ cấp nghề, Trung cấp nghề, Cao đẳng nghề); Đề án 1956
};

% Node 6
\node (n6) [box, below=of n5] {
    \textbf{Giai đoạn 2014 -- nay: Hợp nhất toàn diện theo Luật Giáo dục nghề nghiệp 2014}\\
    \textit{Cơ quan chủ quản:} Tổng cục Giáo dục nghề nghiệp (Bộ Lao động -- Thương binh và Xã hội)\\
    \textit{Mô hình:} Hệ thống giáo dục nghề nghiệp mở, liên thông; Khung trình độ quốc gia (VQF)
};

% Lines with arrows
\draw [arrow] (n1) -- (n2);
\draw [arrow] (n2) -- (n3);
\draw [arrow] (n3) -- (n4);
\draw [arrow] (n4) -- (n5);
\draw [arrow] (n5) -- (n6);

\end{tikzpicture}%
}
\caption{Tiến trình lịch sử và sự chuyển biến thể chế quản lý hệ thống giáo dục nghề nghiệp Việt Nam (1945 -- nay)}
\label{fig:so_do_tien_trinh}
\end{figure}
